\subsection{Parentheses, Brackets, and Quoting}

Parentheses normally express a Lisp function call:
\code{(+ 2 3)} calls addition. Brackets have several roles.
A recognized operator makes them an ABCL form, such as object creation,
assignment, or message transmission. Otherwise, they construct a list
of evaluated values. Thus \code{[:add 2 3]} constructs the data
\code{(:ADD 2 3)}; it does not call an addition method by itself.
The receiver's script gives that data meaning.

\begin{lstlisting}
(let ((x 7))
  [:point x (+ x 1)])     ; => (:POINT 7 8)
'(red green blue)        ; Literal symbols, not variable lookups
[:label "red"]           ; A keyword and a string
\end{lstlisting}

A keyword such as \code{:point} evaluates to itself. An ordinary symbol
such as \code{x} denotes a variable in an expression, but becomes a
binding variable in a pattern. The same bracket syntax therefore serves
as construction at the sending site and decomposition at the receiving site.
Semicolons introduce comments to the end of a line.

\subsection{Matching Data Inside a Script}

Use \code{match} when an already selected handler needs to inspect a value.
Clauses use \code{is}, and a final \code{otherwise} clause provides a fallback.
The wildcard \code{\_} matches a value without introducing a named binding.
A dotted pattern separates a list's head from its remaining tail.

\sourcelisting[firstline=14,lastline=21]{tutorial/examples/01-basics.abclr2}

\begin{lstlisting}
[classifier <== [:classify [:point 2 3]]]
; => (:POINT 2 3)
[classifier <== [:classify [1 2 3]]]
; => (:NONEMPTY 1 (2 3))
[classifier <== [:classify nil]]   ; => :EMPTY
[classifier <== [:classify 9]]     ; => :ATOM
\end{lstlisting}

The point case must precede the broad head/tail case in this \code{match}:
a point is also a nonempty list. Pattern matching tests structure;
it does not validate that coordinates are numbers. Use \code{if} in
the selected body for such a check.

\subsection{Temporary Variables and Loops}

An object's \code{state} survives successive messages.
A \code{temporary} declaration at the beginning of a handler instead
creates local variables for that invocation. Here the accumulator starts
at zero every time:

\sourcelisting[firstline=23,lastline=30]{tutorial/examples/01-basics.abclr2}

\begin{lstlisting}
[list-sum <== [:sum [2 3 4]]]   ; => 9
[list-sum <== [:sum nil]]       ; => 0
\end{lstlisting}

The \code{do} loop binds \code{rest}, tests for termination, runs its body,
and updates \code{rest} using \code{cdr}. You can also use \code{let}
for a local accumulator. This distribution provides \code{dotimes}
and \code{dolist} macros that expand into supported \code{do} forms.
Do not assume that every host Lisp macro expands into forms the ABCL
parser understands.

\paragraph{Exercise 1.}
Define \code{sum-of-squares} with a \code{[:sum values]} handler. Reuse the loop structure,
but add the square of each number. Check that an empty list returns zero.
Compare your reasoning with Section~\ref{sol:squares}.
